\section{Ограничения и насоки за бъдещо развитие}

Всяка инженерна разработка се реализира в определени граници: времеви, технологични, организационни и методологични. Зрелият технически анализ не се изчерпва с описване на постигнатото, а включва и ясно формулиране на онова, което съзнателно или фактически е останало извън обхвата. Настоящият проект не прави изключение. Макар да реализира завършен учебно-приложен модел на онлайн магазин, в него съществуват редица ограничения, които очертават естествена посока за развитие към по-стабилна, сигурна и мащабируема система.

Тази глава систематизира именно тези ограничения и ги превежда в конкретни насоки за бъдеща работа.

\subsection{Ограничения на текущата реализация}

\subsubsection{Опростен модел на сигурност}

Най-значимото ограничение на системата е свързано със сигурността. Макар паролите да се съхраняват чрез \texttt{Pbkdf2PasswordEncoder}, достъпът до защитените страници и операции се контролира предимно чрез ръчни проверки върху \texttt{HttpSession}. Това решение е разбираемо в учебен контекст, но не осигурява достатъчно високо ниво на централизираност, проследимост и устойчивост на политиката за достъп.

Липсват роля-базирана авторизация, security filter chain, формален security context и по-фина диференциация между различни типове потребители. В production среда това би ограничило възможността да се отделят административни, оперативни и клиентски действия.

\subsubsection{Опростен модел на поръчките}

Втората важна граница на проекта е домейнното моделиране на поръчките. Текущата реализация използва \texttt{HistoryEntity} като обобщение на завършена покупка, но не поддържа отделни позиции, статуси, платежни състояния, номер на доставка или жизнен цикъл на поръчката. Това означава, че системата е способна да архивира резултата от checkout операцията, но не и да управлява богат order management процес.

За целите на дипломната разработка тази опростеност е приемлива и дори полезна, защото поддържа домейн модела в разумни граници. В същото време тя ясно показва какво би следвало да бъде добавено при по-сериозно развитие на приложението.

\subsubsection{Непълна изолация на данните по потребител}

Част от логиката за количката и продуктите работи върху глобални колекции, вместо да бъде строго филтрирана по текущия потребител. Това е едно от най-важните технически ограничения на текущия код. В еднопотребителски или демонстрационен контекст проблемът може да остане скрит, но в реална многопотребителска среда той би довел до логически несъответствия.

Този тип ограничение има висока приоритетност, защото засяга не само архитектурната чистота, но и коректността на бизнес поведението.

\subsubsection{Конфигурация с чувствителни данни}

Наличието на credentials в \texttt{application.properties} е operational компромис, който е често срещан в ранни фази на разработка, но не е приемлив за production практика. Потребителски имена, пароли, mail настройки и локални URL адреси следва да бъдат изнесени в защитена конфигурация, външни environment variables или секретно хранилище.

Макар този проблем да не пречи на демонстрационното стартиране, той ограничава възможността проектът да бъде безопасно споделен, публикуван или разгръщан без допълнителна подготовка.

\subsubsection{Непълна front-end централизация}

Интерфейсът е функционален, но front-end организацията остава сравнително фрагментирана. Навигационната лента е повторена в множество шаблони, CSS често е inline, а в различни страници се използват различни версии на външни библиотеки. Това не нарушава работоспособността, но затруднява по-систематичната поддръжка на визуалния слой.

Отделно от това локализацията е само частично последователна. Част от текстовете са изнесени в \texttt{messages} файловете, но други остават твърдо записани в шаблоните. Това затруднява превръщането на приложението в реално двуезична система.

\subsubsection{Ограничено покритие от автоматизирани тестове}

Макар проектът да използва стандартната Maven структура, включваща и тестова зависимост, липсва развита стратегия за unit, integration или UI тестове. Това означава, че качеството на системата в момента зависи в по-голяма степен от ръчна проверка, кодов анализ и демонстрационни сценарии.

В образователен контекст подобна ситуация е често срещана, но за следващ етап на развитие би било силно препоръчително да се въведе автоматизирана тестова рамка.

\subsection{Приоритетни насоки за бъдещо развитие}

След идентифициране на ограниченията следва логичният въпрос кои подобрения имат най-висока стойност. В контекста на разглеждания проект бъдещото развитие може да бъде структурирано по няколко основни направления.

\subsubsection{Внедряване на Spring Security}

Най-естествената следваща стъпка е пълноценно въвеждане на Spring Security. Това би позволило:

\begin{itemize}[leftmargin=*]
    \item централизиран контрол върху достъпа до маршрути;
    \item разграничаване на роли като администратор, оператор и обикновен клиент;
    \item използване на стандартизиран authentication flow;
    \item по-лесно управление на сесии, logout и защита от неоторизиран достъп.
\end{itemize}

Тази стъпка не само би повишила сигурността, но и би изчистила част от логиката в контролерите, която в момента е отговорна за ръчни проверки на достъпа.

\subsubsection{Рефакториране на количката и поръчковия модел}

Втората високоприоритетна посока е преработване на модела за количка и поръчки. Добро решение би било въвеждане на отделни класове от тип \texttt{Order} и \texttt{OrderItem}, както и по-ясно разграничение между активна количка и историческа поръчка. Така системата би могла да поддържа:

\begin{itemize}[leftmargin=*]
    \item множество позиции в рамките на една поръчка;
    \item статуси на поръчката;
    \item запазване на артикулите в историята без загуба на детайлна информация;
    \item по-прецизна отчетност и възможност за последваща статистика.
\end{itemize}

Паралелно с това следва да бъде преработена логиката на количката така, че всички операции да бъдат строго филтрирани по текущ потребител и по идентификатор на продукта.

\subsubsection{Външна конфигурация и профили за среди}

Operational зрелостта на проекта може значително да се подобри чрез използване на профили за \texttt{development}, \texttt{test} и \texttt{production} среда. Това би позволило разделяне на локални mail настройки, тестова база данни и production конфигурация. Паралелно с това чувствителните данни следва да бъдат изнесени в environment variables или друг сигурен механизъм.

Допълнително би било полезно да се въведат миграции на схемата чрез инструмент като Flyway или Liquibase. По този начин структурата на базата ще бъде version-controlled по по-професионален начин.

\subsubsection{Централизация и модернизация на front-end слоя}

Следваща естествена стъпка е подобряване на организацията на интерфейса. Това може да стане чрез:

\begin{itemize}[leftmargin=*]
    \item извличане на повторяемите навигационни блокове в Thymeleaf fragments;
    \item централизиране на inline CSS във външни стилови файлове;
    \item уеднаквяване на използваните версии на Bootstrap и външни библиотеки;
    \item разширяване и довършване на локализацията;
    \item подобряване на responsive поведението и достъпността на формите.
\end{itemize}

Тези стъпки няма да променят фундаментално бизнес логиката, но значително ще повишат професионалното качество на потребителския интерфейс.

\subsubsection{Автоматизирано тестване}

Въвеждането на тестове е друга важна посока. Една разумна тестова стратегия може да включва:

\begin{itemize}[leftmargin=*]
    \item unit тестове за service логиката;
    \item integration тестове за repository и контролери;
    \item сценарийни тестове за регистрация, вход, checkout и история;
    \item тестове за валидационните правила в binding моделите.
\end{itemize}

Добавянето на тестове би повишило увереността в бъдещи промени и би направило проекта по-подходящ за по-сериозно надграждане.

\subsection{Средносрочни и дългосрочни възможности}

Освен приоритетните подобрения съществуват и по-широки направления за развитие, които биха трансформирали системата от учебен прототип към по-зряло приложение.

\subsubsection{Разширяване на каталога и потребителските функции}

В бъдеще могат да бъдат добавени търсене, сортиране, странициране, по-богато филтриране, маркиране на любими продукти, управление на снимки чрез реално качване на файлове и по-детайлни продуктови описания. Това би направило каталога значително по-полезен и реалистичен.

\subsubsection{Интеграция с външни услуги}

Потенциално развитие е и интеграцията с платежни оператори, куриерски услуги, SMS или push известия, както и по-богат REST API за външни клиенти. Такива стъпки биха повишили приложната стойност на проекта и биха го доближили до реални бизнес системи.

\subsubsection{Контейнеризация и deployment автоматизация}

За production ориентирано развитие логична стъпка е контейнеризация с Docker и по-късно автоматизирано разгръщане чрез CI/CD процес. Това би улеснило повторяемото стартиране на приложението и би намалило зависимостта от конкретната локална среда на разработчика.

\subsection{Предложена пътна карта за надграждане}

За по-голяма яснота насоките за развитие могат да бъдат обобщени в кратка приоритетна пътна карта.

\begin{table}[H]
    \centering
    \caption{Предложена пътна карта за бъдещо развитие}
    \label{tab:future-roadmap}
    \begin{tabularx}{\textwidth}{L{2.5cm}L{4cm}Y}
        \toprule
        Приоритет & Направление & Очакван ефект \\
        \midrule
        Висок & Spring Security и роли & По-сигурен и по-чист модел на достъп \\
        Висок & Рефакториране на количка и поръчки & Коректност в многопотребителска среда и по-богат домейн модел \\
        Висок & Изнасяне на secret конфигурация & По-добра operational безопасност \\
        Среден & Централизация на front-end и локализация & По-добра поддръжка и UX последователност \\
        Среден & Автоматизирани тестове & По-висока надеждност при развитие на кода \\
        Среден & Миграции и профили за среди & По-предсказуем build и deployment процес \\
        Нисък към среден & Разширени търговски функции & По-голяма практическа стойност на приложението \\
        Нисък към среден & Контейнеризация и CI/CD & По-лесно разгръщане и възпроизводимост \\
        \bottomrule
    \end{tabularx}
\end{table}

Таблица \ref{tab:future-roadmap} показва, че развитието на проекта може да бъде планирано стъпково. Най-напред следва да се адресират качествата, които пряко влияят на коректността и сигурността, а след това да се премине към по-богати функционални и operational подобрения.

\subsection{Обобщение}

Ограниченията на разглежданата система са естествени за проект с учебно-приложен характер. Те не намаляват стойността на вече постигнатото, а по-скоро определят границите, в които то следва да бъде интерпретирано. Най-важното е, че почти всички установени ограничения могат да бъдат превърнати в ясни и технически осъществими следващи стъпки.

Точно това превръща проекта в добра основа за бъдещо надграждане. Той не е крайна точка, а стабилен първи слой от по-широка система, която при следваща итерация би могла да стане по-сигурна, по-пълнофункционална и по-близка до production стандартите. Следващата и последна глава обобщава постигнатото и формулира заключителните изводи на дипломната работа.
